% AUTORIDAD DE FUENTE:
% expose_VIII_body.tex, líneas 2103--2170; índices de escaneo 307--310;
% folios de fuente 296--299. Las correcciones materiales de la capa
% francesa se conservan reversiblemente mediante \sisi. El diagrama
% reproduce la dirección de flecha corregida a mano en la fuente.

\medskip
\noindent\underline{Observaciones 6.8.}\quad a) Aplicaremos 6.6 y 6.7
cuando $S$ sea un rasgo y $D$ el divisor definido por la inclusión del
punto cerrado $s$ de $S$ en $S$. En este caso, $U$ es el espectro del
cuerpo de fracciones $K$ de $S$, y $\overline G$ es el modelo de Néron
usual del grupo multiplicativo sobre $K$, introducido por Raynaud.
Obsérvese que, incluso en este caso particular, la demostración de 6.6
y 6.7 utiliza 6.5 sobre bases más generales que los rasgos (como $P$ y
sus potencias cartesianas). Sin embargo, en este caso basta conocer la
\underline{existencia} del modelo de Néron $\overline G$ en la
situación especial considerada: 6.5 para esquemas base lisos sobre
$S$, y por tanto 6.6 y 6.7, se deducen entonces formalmente. Para
nosotros, la introducción de $\overline G$ servirá sobre todo de etapa
intermedia cómoda para estudiar, en la sección siguiente, el problema
de prolongar extensiones y biextensiones por el grupo multiplicativo.

\smallskip
\noindent b) Puede definirse un modelo de Néron--Raynaud $\overline G$
de $\underline G_{mU}$, que satisfaga la condición fundamental (6.4.2)
para todo $S'$ liso sobre $S$, bajo condiciones mucho más generales
que las consideradas anteriormente, siendo $U$ meramente un abierto
esquemáticamente denso de un esquema normal localmente noetheriano
$S$. En este caso $\overline G$ es un \underline{haz} fppf, no
necesariamente representable. Denotemos por $\underline D$ el haz, en
el sitio étale pequeño de $S$, de los divisores sobre $S$ cuyo soporte
está contenido

%%%%%%%%%%  scan idx 308  |  folio 297  %%%%%%%%%%

en $T=S\setminus U$, y, por abuso de lenguaje, denotemos también con
la misma letra su extensión canónica a un haz sobre
$(\mathrm{Sch})_{/S}$. Sus secciones sobre un esquema relativo
$f:S'\to S$ son las secciones de
$f_{\mathrm{\acute{e}t}}^*(\underline D)$ sobre $S'$. El modelo buscado
es una extensión de grupos
\begin{equation}\tag{6.8.1}
 0\longrightarrow\underline G_{mS}\longrightarrow\overline G
 \longrightarrow\underline D\longrightarrow0.
\end{equation}
La construcción de una extensión canónica (6.8.1), sin hipótesis de
normalidad sobre $S$, no presenta dificultad si se usa la descripción
general VII 1.1.6; al mismo tiempo se obtiene (6.4.2) para $S'$ étale
sobre $S$. Para que la misma relación valga para todo $S'$ liso sobre
$S$, equivale a comprobar que la formación del haz $\underline D$
conmuta con todo cambio de base liso $f:S'\to S$. Cuando $S$ es normal,
esto se deduce de EGA $\mathrm{Err}_{\mathrm{IV}}$ 53, 21.4.9.

\medskip
\noindent 6.9.\quad ``\underline{Autocrítica}'' (observaciones de
P.~Deligne).\footnote{El redactor de la presente exposición, habiéndose
dado por vencido, invita al lector, si fuera necesario, a reescribir
las secciones 5 y 6 siguiendo las indicaciones adjuntas.}

Las secciones 5 y 6 son realmente poco inspiradas. Periódicamente se
tiene la impresión de que se rehace ``a mano'' el formalismo de
$i_*,i^*$, sin distinguir lo que es geométrico de lo que es
``general nonsense''.

Propongo extraer los enunciados ``geométricos'' a), b), c):

\smallskip
\noindent a) (cf. 6.8 b)) Sea
$j:\sisi{V}{U}\hookrightarrow S$ un abierto esquemáticamente denso de
un esquema normal noetheriano $S$. El haz étale $\underline G_{mS}$ se
identifica entonces con un subhaz del haz étale
$j_*\underline G_{mU}$. Sea $\underline D$ el haz étale cociente en la
sucesión exacta
\[
 0\longrightarrow\underline G_{mS}\longrightarrow
 j_*\sisi{\underline G_{mS}}{\underline G_{mU}}
 \longrightarrow\underline D\longrightarrow0.
\]

%%%%%%%%%%  scan idx 309  |  folio 298  %%%%%%%%%%

\underline{Entonces} la formación de $\underline D$ conmuta con todo
cambio de base liso (EGA $\mathrm{Err}_{\mathrm{IV}}$ 53, 21.4.9);
\underline{además}, $R^1j_*\underline G_m=0$ (Teorema 90) si los
anillos estrictamente locales de $S$ son factoriales.

\smallskip
\noindent b) Si $U=S\setminus Y$, donde $Y$ es un divisor
geométricamente unirrama, y si $i$ es la inclusión de $Y$ en $S$,
entonces
\[
 \underline D=i_*\mathbb Z_Y.
\]

\smallskip
% Canon: SGA-SRC-001273.
\noindent c) \sisi{Si $Y$ es un subesquema cerrado de un esquema $X$ y
$E$ es un conjunto}{Si $Y$ es un subesquema cerrado de un esquema $X$,
si $i:Y\hookrightarrow X$ designa la inmersión cerrada, y si $E$ es un
conjunto}, entonces el haz étale $i_*\underline E_Y$ está
representado por un esquema étale sobre $X$, cuya formación conmuta con
todo cambio de base.

\medskip
\noindent\underline{Corolario.} \underline{Sea $S$ un esquema normal
noetheriano, sea $Y$ un divisor geométricamente unirrama y sea
$U=S\setminus Y$:}
\begin{center}
\begin{tikzcd}[column sep=huge]
 U \arrow[r, hook, "j"'] & S & Y \arrow[l, "i"']
\end{tikzcd}
\qquad .
\end{center}
\underline{Existe un único esquema en grupos liso $\mathcal G_m$ sobre
$S$ tal que, después de todo cambio de base liso, se tiene
$\mathcal G_m|S_{\mathrm{\acute{e}t}}=j_*(\underline G_m)_U$.}

En efecto, las condiciones precedentes determinan $\mathcal G_m$ como
haz, para la topología étale, sobre la categoría de esquemas lisos
sobre $S$. Por a) y b), la sucesión
\[
 0\longrightarrow\underline G_m\longrightarrow\mathcal G_m
 \xrightarrow{\ q\ }i_*\mathbb Z_Y\longrightarrow0
\]
es exacta. Por c), $i_*\mathbb Z_Y$ es representable. Finalmente, el
morfismo $q$ es representable.

El uso del ``sitio liso pequeño'' precedente debería permitir deducir
formalmente de a) los enunciados de la sección 6.

%%%%%%%%%%  scan idx 310  |  folio 299  %%%%%%%%%%
